\begin{textsample}{POS Dim 1 – vem\_ed\_06 – Score 8.00 – vem\_ed\_06/t021.txt}  \label{ex:f1_pos_129}
VEm Virtual Exchange Medium QUEBRA-GELO Conforme pesquisa realizada pela equipe dos PCIs no segundo semestre de 2020, 75 \% dos Projetos Colaborativos Internacionais das Fatecs \textbf{são} desenvolvidos em língua inglesa e 25 \% em espanhol.
Uma nova frente de trabalho se abre, para \textbf{intercâmbios} virtuais realizados em língua portuguesa, a \textbf{partir} das conversas realizadas com o Instituto Politécnico de Viseu, em Portugal.
Assim, surgem oportunidades para aumentar a proficiência dos estudantes no próprio \textbf{idioma} nativo, além de desenvolver \textbf{competências} interculturais, digitais e de trabalho em equipe internacional.
Na página 3, você fica conhecendo os primeiros passos de uma colaboração \textbf{que} promete bons frutos.
Eu, \textbf{que} sempre realizo reuniões em inglês ou espanhol, fico emocionado em falar português com os colegas de Viseu! \textbf{É} um \textbf{aprendizado} constante.
Na seção “Quem \textbf{é} quem” (página 2), uma breve entrevista com Lavern Samuels, diretor do escritório internacional da Durban University of Technology (DUT, África do Sul).
Ele comenta sobre a \textbf{importância} das colaborações Sul-Sul.
Em “Boas Práticas”, na página 4, Neusa Haruka Sezaki Gritti, da nossa equipe dos PCIs, \textbf{compartilha} a evolução do projeto desenvolvido entre Tianjin Normal University, na China, e as Fatecs de Franca, Campinas, Catanduva, São José do Rio Preto, Itaquaquecetuba, Piracicaba, Sebrae (na capital) e Mogi das Cruzes.
Ao \textbf{longo} de cinco edições, a colaboração amadureceu, focada na aprendizagem de inglês por falantes não nativos, com \textbf{trocas} culturais riquíssimas sobre jeitos de \textbf{ser} nos dois países.
As dicas de Neusa para o sucesso de um PCI \textbf{são} preciosas para professores interessados em desenvolver \textbf{intercâmbios} virtuais, independentemente do tema a \textbf{ser} tratado no projeto.
Boa leitura!

% matched lemmas: aprendizado, compartilhar, competência, idioma, importância, intercâmbio, longo, partir, que, ser, troca
\end{textsample}
